%%%%%%%%%%%%%%%%%%%%%%%%%%%%%%%%%%%%%%%%%
% The Legrand Orange Book
% LaTeX Template
% Version 2.3 (8/8/17)
%
% This template has been downloaded from:
% http://www.LaTeXTemplates.com
%
% Original author:
% Mathias Legrand (legrand.mathias@gmail.com) with modifications by:
% Vel (vel@latextemplates.com)
%
% License:
% CC BY-NC-SA 3.0 (http://creativecommons.org/licenses/by-nc-sa/3.0/)
%
% Compiling this template:
% This template uses biber for its bibliography and makeindex for its index.
% When you first open the template, compile it from the command line with the 
% commands below to make sure your LaTeX distribution is configured correctly:
%
% 1) pdflatex main
% 2) makeindex main.idx -s StyleInd.ist
% 3) biber main
% 4) pdflatex main x 2
%
% After this, when you wish to update the bibliography/index use the appropriate
% command above and make sure to compile with pdflatex several times 
% afterwards to propagate your changes to the document.
%
% This template also uses a number of packages which may need to be
% updated to the newest versions for the template to compile. It is strongly
% recommended you update your LaTeX distribution if you have any
% compilation errors.
%
% Important note:
% Chapter heading images should have a 2:1 width:height ratio,
% e.g. 920px width and 460px height.
%
%	Cities list:
%	New York
%	Beijing
%	London (city)
%	Singapore
%	Dubai
%	Tokyo
%	Los Angeles
%	Frankfurt
% 	Kuala Lumpur
%	Brussel
%	Milan
%	Berlin
%
%%%%%%%%%%%%%%%%%%%%%%%%%%%%%%%%%%%%%%%%%

%----------------------------------------------------------------------------------------
%	PACKAGES AND OTHER DOCUMENT CONFIGURATIONS
%----------------------------------------------------------------------------------------

\documentclass[11pt,fleqn,oneside]{book} % Default font size and left-justified equations

%----------------------------------------------------------------------------------------

\input{structure} % Insert the commands.tex file which contains the majority of the structure behind the template

\begin{document}

%----------------------------------------------------------------------------------------
%	TITLE PAGE
%----------------------------------------------------------------------------------------

\begingroup
\thispagestyle{empty}
\begin{tikzpicture}[remember picture,overlay]
    \node(background) at (current page.center) {\includegraphics[height=\paperheight]{background}};
\end{tikzpicture}
\endgroup

\newpage

% \begingroup
% \section{Abstract}\index{Abstract}
% \thispagestyle{empty}
% \addcontentsline{toc}{chapter}{\textcolor{trolleygrey}{Abstract}}
\usechapterimagefalse % If you don't want to include a chapter image, use this to toggle images off - it can be enabled later with \usechapterimagetrue
\chapter{Abstract}
\textbf{Ripa Exchange è un exchange ibrido-decentralizzato con il focus di abbassare il livello di ingresso per aprire 
nuove piattaforme di exchange ed al contempo offrire ai traders un ambiente sicuro ed efficiente per eseguire operazioni di 
trading giornaliere.}\\

Il Team di Ripa Exchange crede che, indipendentemente dagli sviluppi nel mondo delle criptomonete, è oggigiorno costoso per aprire,
gestire, e trovare fiducia nelle proprie operazioni di cambio valuta non solo per le risorse richieste per avviare
una piattaforma di scambio affidabile ma soprattutto per sviluppare la piattaforma in sè e trovare la liquidità 
necessaria per generare profitto nei primi 3-5 anni di attività.\\

Azione è richiesta ed azione è richiesta adesso. Gli utenti sono frustrati da exchange inaffidabili che chiudono scappando con i 
fondi, vengono hackerati oppure non sostengono il carico giornaliero di una'industria crescente come questa è.
Nonostante l'impegno degli imprenditori di exchange per offrire piattaforme di cambio valuta efficienti, affidabili e 
semplici da utilizzare i fondi necessari per costruire queste piattaforme sono nella fascia di 250-300 mila euro senza includere
il costo del personale per offrire supporto di livello platino agli utenti finali, l'infrastruttura di gestione, e costi giornalieri
dell'attività. Tali costi vengono sostenuti per avere una piattaforma discreta alla quale ci si lega ad unica via con una software
house per sviluppi futuri.\\

É obbiettivo di questo progetto di sviluppare una piattaforma di cambia valuta Open Source, efficiente, affidabile ed 
offrire liquidità\footnote{Grazie alla tecnologia RLSP (Ripa Liquidity Service Provider)} all'exchange così creato dal \textbf{primo}
giorno di attività in modo tale che gli imprenditori in questo settore possono focalizzarsi nel pubblicizzare la propria piattaforma,
offrire supporto tecnico di libello platino ai propri utenti e soddisfare tutte le eterogenee legislazioni nell'industria delle valute virtuali.
Mentre vogliamo che l'esperienza di trading per l'utente finale sia la migliore possibile offrendo un ambiente di trading sicuro ed affidabile.\\
\usechapterimagetrue
% \endgroup

%----------------------------------------------------------------------------------------
%	TABLE OF CONTENTS
%----------------------------------------------------------------------------------------

\chapterimage{chapter_head_1_Beijing.jpg} % Table of contents heading image
\renewcommand*\contentsname{Indice dei Contenuti}
\tableofcontents % Print the table of contents itself

%----------------------------------------------------------------------------------------
%	PART
%----------------------------------------------------------------------------------------


%\part{Part One}

%----------------------------------------------------------------------------------------
%	CHAPTER 2: Introduction
%----------------------------------------------------------------------------------------

\chapterimage{chapter_head_2_London.jpg} % Chapter heading image

\chapter{Introduzione}
Il livello di accesso all'industria delle valute virtuali è piuttosto elevato da un punto di vista tecnico per l'utente medio, 
ed ha un livello di accesso economicamente elevato per l'imprenditore che vuole avviare un'attività in questo campo, valore dato 
dall'acquisto del codice sorgente dell'exchange di criptovalute, dal reclutamento di personale operativo per le operazioni di DevOps,
per reclutare operatori di supporto tecnico, per sottostare alle leggi nazionali ed internazionali in materia di AML/KYC, per avere liquidità
dal primo giorno di avvio delle operazioni di cambio. Con questo progetto, Ripa Exchange, vuole abbassare questo livello di ingresso 
perchè \textbf{gestire un exchange è DIFFICILE} ed il Team Ripa Exchange vuole che il lettore e futuro imprenditore nel campo delle valute
virtuali si concentri/focalizzi su cose veramente importanti non cavilli burocratici che l'industria richiede perchè TU vuoi avviare un'attività
in questo campo e tu necessiti del codice sorgente per avviarla (ed il codice sorgente per avviarla ti deve essere concesso gratuitamente).

Per rafforzare/chiarire questo principio economico il punto è che per avviare un exchange è richiesto un investimento esiguo dal tuo 
venture capital di riferimento e anche con un tale investimento il ritorno dei fondi investiti all'angel investor non sono garantiti
nei primi 5 anni di attività.

Per avviare un servizio professionale di cambio valuta crediamo che il codice sorgente dell'exchange e la liquidità da offrire ai tuoi clienti
dal primo giorno di attività dovrebbero esserti forniti gratuitamente: non deve essere accettabile pagare \euro150.000,00 ad una software house
solo per avere una piattaforma che funzioni e per la quale dovete erogare altri \euro100.000,00-150.000,00 per brandizzarla, personalizzarla
per i vostri scopi legandovi in questo modo ad una compagnia/software house che potrebbe fallire in futuro e che vi farebbe ritrovare
senza poter fare ulteriori modifiche al vostro exchange in quanto non vi hanno mai consegnato il codice sorgente della piattaforma/del prodotto
su cui il vostro business fa affidamento/si basa.

\textbf{Noi crediamo che tutto questo debba essere gratuito}: dobbiamo offrirti la miglior tecnologia nel mercato attuale in modo tale che tu 
possa focalizzarti sulla tua attività mentre noi ci focalizziamo in sviluppare la tecnologia per gestire la tua attività/i tuoi affari
in maniera efficiente, sicura, responsiva e produttiva. Questo è il motivo per cui Ripa Exchange si prefigge lo scopo di sviluppare una rete di exchange
focalizzandosi/utilizzando un'architettura che è \textit{efficiente, sicura, responsiva, compliant e personalizzabile} in modo tale che ogni exchange
nella rete può basarsi su di fondamenta solide mentre è comunque possibile modificare la singola istanza secondo le personalizzazioni 
che la compagnia che gestisce quell'istanza richiede/quel singolo exchange richiede.

Per raggiungere questo scopo abbiamo scelto di sviluppare la tecnologia Ripa Liquidity Service Provider basata su ARK - a blockchain for consumer adoption - 
il cui obbiettivo primario è incrementare l'adozione delle tecnologie blockchain agli utenti finali focalizzandosi su due aree critiche: 
\underline{Una Efficiente e Sicura Tecnologia di Base} e \underline{Servizi Reali per Gente Comune}. L'ecosistema ARK è ancora ad uno stadio
iniziale di sviluppo: la futura implementazione di ARK 2.0 permetterà di eseguire smart contract nativi sulla blockchain, che permetteranno
alla blockchain ARK di competere con la sua concorrente Ethereum da un punto di vista tecnologico. RLSP permetterà a tutti gli exchange della rete Ripa
di condividere la stessa liquidità-lo stesso libro degli ordini in base ai parametri di condivisione della singola istanza dell'exchange Ripa
configurabili nella tua area amministratore.

Il Ripa Founder Team (RFT), come presentato in ripaex.io, opera nel nome del gruppo Ripa. L'RFT è responsaile per l'uso corretto dei finanziamenti
ottenuti nella campagna di scambio token RIPA TEC come spiegato di seguito in questo documento.

Il RFT comprende-impone che i risultati delle fasi di Prevendita e RIPA TEC saranno usati esclusivamente per il finanziamento del progetto 
\emph{Ripa Exchange} come spiegato in questo whitepaper, che saranno resi disponibili per estrazione nella piattaforma apposita resta disponibile
più tardi quest'anno tec.ripaex.io, e che dovranno risultare nella creazione di una entità legale che sarà chiamata \emph{Ripa Exchange}. La creazione di tale 
compagnia è pianificata per il primo quadrimestre del 2019.

Fino a tale data RFT opera invece di \emph{Ripa Exchange, una compagnia nel processo di essere incorporata}.


%------------------------------------------------

%This statement requires citation \cite{article_key}; this one is more specific \cite[162]{book_key}.

%------------------------------------------------

\section{Terminologia}\index{Terminologia}
\begin{description}
	\item[\textsc{Ripa Exchange:}] un exchange FIAT <-> CRYPTO sviluppato partendo dal codice sorgente di Peatio \cite{peatio}
	\item[\textsc{Ripa Blockchain:}] una blockchain DPOS in cui la liquidità ed il libro degli ordini sono condiviss tra tutti gli exchange della rete
	\item[\textsc{Ripa Token - XPX:}] un token crittograficamente sicuro tramite il protocollo DPOS scambiato sulla blockchain Ripa 
	\item[\textsc{RIPA:}] l'ecosistema finanziario DPOS composto da Ripa Exchange e Ripa Blockchain
	\item[\textsc{RIPAEX:}] il nome del progetto, sito web del progetto e domini del progetto
	\item[\textsc{RLSP:}] Ripa Liquidity Service Provider, un libro ordini condiviso tra gli exchange nella stessa rete Ripa
	\item[\textsc{RipaEx ICO:}] il nome del periodo di cambia valuta a tasso fisso composto dalle fasi di PreSale e RIPA TEC	
	\item[\textsc{ARK:}] una piattaforma per l'adozione consumer di tecnologie blockchain \cite{ark}
	\item[\textsc{ACES:}] Ark Contract Execution Services \cite{aces} fornisce semplici protocolli e strumenti per costruire robusti
	mercati di scambio tra blockchain basati sulla tecnologia ARK SmartBridge
	\item[\textsc{“,” or “.”:}] la notazione italiana per virgola decimale e punti di migliaia è stata scelta per la stesura di questo 
	whitepaper: detto questo “,” rappresenta un punto decimale e “.” rappresenta la separazione tra migliaia e multipli di migliaia 
\end{description}

%Lists are useful to present information in a concise and/or ordered way\footnote{Footnote example...}.
\pagebreak
\section{Roadmap}
Saranno percorse essenzialmente quattro fasi del progetto RipaEx:\\

\tcbset{roadmapBox/.style={colback=yellow!10!white,colframe=azure(colorwheel),
equal height group=nobefaf,width=(\linewidth-1pt), height from=4cm to 8cm, nobeforeafter,
	center title,
	valign=top, halign=left}}
\begin{center}
\begin{tcolorbox}[roadmapBox,
	title=\textbf{\textsc{Finanziamneto del progetto: XPX presale e RIPA TEC (WP2)}}]

	Questa fase riconosce l'esistenza di un interesse in questo sviluppo di mercato
	da un capo del Mondo all'altro in riguardo l'abbassamento del livello economico di ingresso
	per sviluppare exchange di cryptovalute.
	L'obbiettivo è di eseguire la prima analisi omnicomprensiva dello stato dell'arte per fondare
	le basi delle fasi sucessive del progetto e per sviluppare il primo prototipo
	funzionante di un exchange centralizzato con il codice di Peatio.\\
	\vspace{1cm}
	\centering\textbf{\textsc{Prevista fine della fase: Gennaio 2019}}.
\end{tcolorbox}

\resizebox{0.05\textwidth}{26pt}{$\Downarrow$}

\begin{tcolorbox}[roadmapBox,
	title=\textbf{\textsc{Apertura del primo exchange e sviluppo di strumenti e risorse (WP3)}}]

	La seconda fase prende i risultati della prima e sviluppa da qui
	una serie di strumenti e risorse che offrono guida concisa ed omnicomprensiva per 
	i player senza considerare lo Stato di incorporazione/sviluppo del progetto. Con la prima
	istanza di Ripa Exchange in esecuzione/aperta al pubblico i primi contatti con altri player economici
	nell'industria delle valute virtuali possono essere presi.\\
	\vspace{1cm}
	\centering\textbf{\textsc{Prevista fine della fase: Giugno 2019}}.
\end{tcolorbox}

\resizebox{0.05\textwidth}{26pt}{$\Downarrow$}

\begin{tcolorbox}[roadmapBox,
	title=\textbf{\textsc{Esecuzione, diffusione (WP 7/8) e Coordinamento del Progetto (WP1)}}]

	Durante tutta la durata del progetto,
	attività di diffusione (WP 7/8) sono eseguite nelle quali i risultati dei singoli
	pacchetti di lavoro individuali sono consegnati ai relativi gruppi obbiettivo inclusi
	partners del progetto, supporters RipaEx, manager di exchange, partner bancari come anche
	tutti gli altri gruppi obbiettivo rilevanti al progetto RipaEx/in esecuzione.
	Questa fase copre una vasta scala di tecniche di diffusione, da libri stampati ed 
	elettronici, a workshop e sessioni di training/hackatons, reti esistenti, tutti aventi 
	come obbiettivo di definire standard per la comunicazione tra exchange tra entità pubbliche e private.
	Un OVERARCHING pacchetto di lavoro riguardante la gestione del progetto dall'inizio alla fine, 
	assicurando coordinazione adeguata, quality assurance e controllo del budget (WP1).
\end{tcolorbox}

\resizebox{0.05\textwidth}{26pt}{$\Downarrow$}

\begin{tcolorbox}[roadmapBox,
	title=\textbf{\textsc{Sviluppo di un exchange ibrido-decentralizzato (WP 4-6)}}]

	Usando gli strumenti e le risorse sviluppate in WP3, il pacchetto di lavoro 4-6
	si focalizza mettere in pratica la conoscenza acquisita e gli strumenti. I tre
	pacchetti di lavoro riflettono tre punti focali (e gruppi target) nel network di 
	exchange creato per incorporare/stabilire demo di successo su scala locale, nazionale ed internazionale:
	l'incorporazione di istanze di Ripa Exchange locali (WP4), analisi tecnica della tecnologia
	Ripa Liquidity Service Provider (WP5), e primo MVP/POC dell'exchange ibrido decentralizzato (WP6).
	La fase demo dal cuore dell'azzione RipaEx; WP 2 e 3 focalizzano in produrre deliverables (consegne ad esempio
	strumenti) che abilitano attività di dimostrazione efficienti e di successo.\\
	\vspace{1cm}
	\centering\textbf{\textsc{Prevista fine della fase: Gennaio 2021}}.
\end{tcolorbox}
\end{center}

\section{RipaEx Partners - RipaEx Governance}\index{RipaEx Partners - RipaEx Governance}
La maggior parte dei parnter sono imprenditori nell'industria delle valute virtuali, tuttavia
un Istituto di Ricerca ed Organizzazioni Finanziari sono anche presenti. I partner sono:
\begin{description}
	\item[Coordinatore]: Ripa Exchanges Ltd
	\item[CoBeneficiari]: 
\end{description}
\subsection{RipaEx Governance}\index{RipaEx Governance}
La governance per la rete di exchange creati, lo sviluppo del codice sorgente, il possesso dei token XPX, la Fondazione Ripa.

\section{Sommario del Progetto}
\begin{enumerate}
	\item \textsc{cosa}: RipaEx è un progetto per facilitare la presa di standard per condividere liquidità tra mercati di asset
	criptosicuri (crypto asset marketplace).
	L'obbiettio di RipaEx è la promozione di codice sorgente condiviso per portafogli ed exchange nell'industria delle valute
	virtuali: è obbiettivo di questo documento di riferimento di dare informazioni dettagliate per futuri/prospective,
	sviluppatori di exchange oppure imprenditori in questo campo, per abilitare processi di decision-making corretti e per 
	assicurare il successo dei loro progetti proposti.
	IT SEEKS di analizzare il potenziale reale di applicazione nel Paese scelto di una rete di exchange e la sua posizione nel mercato.
	\item \textsc{cosa}: gli asset crittograficamente sicuri (aka: "crypto assets") sono un'alternativa agli asset emessi dalle borse
	valori dei vari Paesi che ne possiedono una. Anche se, tuttavia, molte borse valori mondiali danno la possibilità ai loro 
	utenti di verificare e gestire gli asset in loro possesso il processo non è sempre trasparente e semplice da fare e questo è il motivo 
	per cui a partire dal 2009 \cite{bitcoin} un nuovo tipo di asset (verificabile dalla comunità) è stato implementato
	per dare a piccoli, medi e grandi investitori completa trasparenza nella gestione dei loro asset di investimento.
	\item \textsc{quanto}: recenti sviluppo a livello dell'Unione Europea e globalmente stanno trasformando sia come le valute virtuali
	sono trattate sia il modo in cui le ICO (Initial Coin Offering) sono legislate. Questi sviluppi combinati hanno messo
	l'uso e la produzione di valute virtuali in una luce favorevole in costante ascesa. \\
	In Ottobre 2015 la Corte Europea di Giustiza ha sentenziato che bitcoin ed altre valute virtuali sono esenti dalla tassazione IVA. \\
	In Luglio 2016 la Commissione Europea ha adottato proposal per ammendare la legislazione vigente 4th Anti-Money Laundering Directive (4AMLD) 
	che introduca le figure degli exchange di valute virtuali ed offritori di wallet all'interno del framework europeo di adeguata verifica
	dell'utente ed antiriciclaggio \cite{EUAMLCrypto}.\\
	In Febbraio 2018 la Commissione Europea lancia l'Osservatorio EU sulle Blockchain (EU Blockchain Observatory and Forum) \cite{EUBOaF} 
	per osservare gli sviluppi chiave nella tecnologia blockchain, promuovere attori europei e rafforzare l'ingaggio europeo con
	diversi stakeholders coinvolti in attività correlate alla blockchain.
	\item \textsc{perchè}: c'è ancora poca regolamentazione sulle ICO e solo gli Stati Uniti d'America al momento
	hanno scritto una legislazione definendo i token delle ICO come securities \cite{SECICO}.
	\item \textsc{chi}: i risultati delle indagini di mercato indicano che il medio tech savvy dai 18 ai 45 anni di età è
	l'utente medio delle valute virtuali tuttavia compagnie in finanza stanno anche comincianto ad inserire schemi di 
	valute virtuali all'interno dei loro portafogli specialmente dalla presentazione del contratto bitcoin futures 
	da parte del Gruppo CME Group Inc. nella borsa valori di Chicago lo scorso 18 Dicembre 2017.
	\item \textsc{quanto}: la capitalizzazione di mercato totale dell'industria delle valute virtuali è stata stimata attorno
	ai 317 miliardi di USD \footnote{Dati Coinmarketcap Aprile 2018}
	ed è prevista la crecita a 5.000 miliardi di USD nel prossimo timeframe di 10 anni \cite{cryptoMCTenYears}.
	\item \textsc{dove}: autorità locali stanno lavorando con i Governi Nazionali per essere certi che gli exchangers sul 
	territorio nazionale di riferimento seguono le procedure nazionali ed internazionali di AML/KYC.
	Venture Capital ed Angel Investors stanno cominciando a rilasciare soluzioni di finanziamento a start-up nell'industria delle
	valute virtuali - FinTech in tutto il mondo, dall'America all'Asia passando per l'Europa ed alcuni Paesi
	stanno avviando schemi di valuta virtuale statali per testare lo scambio di beni e servizi sulle tecnologie a libro 
	mastro condiviso (DLTs - Distributed Ledger Technologies \cite{petro}.
	\item \textsc{quanto costa}: il costo medio per avviare il proprio crypto assets marketplace è attorno a \$ 150.000,00, al pagamento di tale somma
	avrete unicamente un'istanza in esecuzione della vostra piattaforma di cambio: al suddetto prezzo dovete aggiungere i costi
	di personalizzazione della piattaforma prima del lancio ed in futuro, pubblicizzazione della piattaforma, costo di gestione dei server, 
	amministratori di rete, operatori del centro di supporto e costo del dipartimento legale per seguire le leggi AML/KYC dello Stato
	in cui avete deciso di incorporare il vostro exchange e per gestire l'amministrazione della compagnia dietro di esso.\\
	\textbf{Questi costi nascosti sono la ragione per cui noi di RipaEx riteniamo che possedere il codice sorgente del proprio exchange è il miglior
	modo per avviare e mantenere un'attività in questa industria}.
	\item \textsc{come}: il problema principale nell'eseguire un'operazione di cambio FIAT <-> CRYPTO è di trovare partner bancari affidabili
	e nel seguire le differenti procedure AML/KYC del Paese o dei paesi di incorporazione.
	\item \textsc{cosa}: le operazioni di cambio valuta virtuale classiche sono le seguenti:
		\begin{enumerate}[label*=\arabic*.]
			\item \textbf{scambio ad una-via}: nella quale un'applicazione centralizzata ha tutta la liqudità da offrire ai potenziali utenti
			\item \textbf{scambio a due-vie}: nella quale un exchange centralizzato o decentralizzato accoppia richiesta di vendita con richiesta di aquisto degli utenti
		\end{enumerate}
		Dalla classificazione data sopra una ulteriore sotto-classificazione è d'obbligo:
		\begin{enumerate}[label*=\arabic*.]
			\item \textbf{scambi FIAT <-> CRYPTO}: nella quale lo scambio avviene tra valuta FIAT\footnote{Valute emesse dalle banche centrali: EUR, USD, GBP, JPY, altre...}
			e valuta virtuale
			\item \textbf{scambi CRYPTO <-> CRYPTO}: nella quale lo scambio avviene tra valuta virtuale e valuta virtuale
		\end{enumerate}
	Dalle quattro classificazioni sopra si può costruire una matrice con quattro possibili configurazioni per sviluppare la piattaforma di cambia valuta virtuale
	che preferite.
	\item \textsc{con cosa}: le specifiche da guardare quando scegliere la propria piattaforma di exchange sono le seguenti: 
		\begin{enumerate}[label*=\arabic*.]
			\item \textbf{codice}: Open Source, Closed Source oppure soluzione ibrida
			\item \textbf{modularità}: separazione tra l'engine di cambio (engine di accoppiamento ordini), la UI e l'anagrafica utenti
			\item \textbf{responsività della UI}: la UI deve essere responsiva su tutti i device (desktop e mobile)
			\item \textbf{conformità}: con gli standard industriali correnti
			\item \textbf{personalizzazione}: dell'engine di cambio, delle valute da offrire, della UI e di altri aspetti della piattaforma di cambio
			\item \textbf{sicurezza} dei fondi nell'exchange: possibilità di configurare il salvataggio dei fondi in cold wallet e hot wallet
			\item \textbf{trasparenza} dei fondi: proof of solvency dell'exchange
			\item \textbf{Multi-Accounts trading}: possibilità di configurare facilmente nuovi mercati da offrire agli utenti
			\item \textbf{Multi-Accounts users}: possibilià di loggarsi nella piattaforma tramite Google, Facebook, Twitter, e conformità
			con gli standard della FIDO Alliance per le credenziali personali
		\end{enumerate}	
	\textbf{Queste non sono solo decisioni tecniche da prendere ma anche economiche} specialmente il possesso o meno del codice sorgente della propria 
	piattaforma	crypto asset marketplace per permettere le future personalizzazioni del proprio exchange in maniera indipendente invece di affidarsi
	ad una singola software house che esegua le personalizzazioni per voi.
	\item \textsc{come}: possibili operazioni per trovare nuovi utenti per le vostre operazioni di cambia valuta virtuale sono:
	campagne di marketing "ad hoc", funzionalità innovative nell'industria, livelli di tassazione del trading basati sulla quantità, 
	bonus di registrazione e per quantità di scambiato giornaliero/settimanale/mensile, marketing affiliato per gli utenti che portano
	amici a scambiare nell'exchange.
	\item \textsc{come}: per incorporare una crypto asset marketplace un progetto deve tenere in considerazione la seguente legislazione:
		\begin{enumerate}[label*=\arabic*.]
			\item \textbf{AML/KYC}: \textit{Fourth Anti-Money Laundering Directive} se l'attività viene avviata nell'Unione Europea \cite{4AMLD} 
			oppure la legge AML/KYC di riferimento nel Paese di incorporazione (ad esempio \textit{Intelligence Reform \& Terrorism Prevention Act of 2004}
			emessa dalla FinCEN per incorporazione negli Stati Uniti d'America).\\
			Raccomandazioni internazionali per eseguire verifiche AML/CFT sono data dalla Financial Action Task Force on Money Laundering \cite{FATF}.
			\item \textbf{Licenza di Pagamento}: l'attività decisamente più ardua in riguardo alla legittimazione delle operazioni di cambio FIAT <-> CRYPTO
			è l'ottenimento di una \textit{Licenza PSD} \cite{PSD}.
			La licenza PSD seguen la Direttiva del Concilio Europeo 2007/64/EC ed è applicata in ogni Paese attraverso le sue leggi nazionali.
			Il costo di una licenza PSD può variare da \euro XXXX ed \euro XXXX, in riguardo alla volume di affari.
		\end{enumerate}
	\item \textsc{chi}: la natura dell'attività in considerazione del progetto Ripa Exchange (una rete di piccoli exchange FIAT <-> CRYPTO)
	implica che ogni exchange debba avere uno staff di 7-8 persone: N.2 sviluppatori, N.1 operatore di rete/sicurezza, 
	N.1 commerciale amministrativo, N.2 operatori di supporto tecnico, N.1 legale e consigliere sulla tassazione. \\
	Il turnover di una tale compagnia tuttavia, considerato l'alto valore aggiunto del prodotto scambiato, e possibile che raggiunga cifre
	superiori ad \euro 350.000 annui e possibilmente di un'ordine di grandezza superiore. Un'attività di questa dimensione
	suggerisce automaticamente le possibili forme societarie: una ditta individuale, una società a responsabilità limitata,
	una società non-profit o attività sociale, una cooperativa gestita dagli amministratori delegati dei vari exchange.
	Le Agenzie Finanziarie sono attori fondamentali nell'indicare questa scelta, tuttavia la tipologia di attività da
	incorporare dipende dallo status legale che si vuole dare alla stessa e questo varia da Paese a Paese.		
	\item \textsc{come}: una risorsa di fondi per una marketplace di dimensioni piccole-medie può essere: prestito bancario,
	prestito privato a basso interesse, credito commerciale, equity finanziario, venture capital di business angels. Avendo un
	business plan consolidato e garanzie finanziarie adeguate sono elementi chiave nell'assicurarsi i finanziamenti necessari
	allo sviluppo che si vuole eseguire. Il Fondo Europeo di Investimento (EIF) dell'EIB offre supporto nella forma di 
	garanzia per le SME (piccole e medie imprese).
	\item \textsc{chi}: gli argomenti per l'incorporazione di un crypto asset marketplace sono la libertà finanziaria, la
	decentralizzazione delle operazioni di trasferimento del valore ed il reale possesso del proprio denaro.
	Ci sono altri benefici, ben documentati, come pagamenti veloci, guadagno a lungo termine dato da un'economia deflazionaria, previsione
	delle depressioni economiche come quella del 2008, ma soprattutto le valute virtuali sono l'unico competitor diretto alle operazioni
	di trasferimento del valore centralizzate eseguite dalle banche centrali. d value-transferring operations done by central banks.
	\item \textsc{perchè}: c'è consenso in letteratura che l'uso delle valute virtuali invece delle valute fiat permette un livello di
	libertà finanziaria maggiore specialmente perchè queste tecnologie di trasferimento valore rientrano perfettamente nella scuola
	di economia Austriaca \cite{austrianTheory} (TODO: aggiungere sulla scuola di economia Austriaca)
	\item \textsc{perchè}: benefici delle tecnologie di valuta virtuale (ToDo: aggiungere numeri)
	\item \textsc{come}: securizzare asset sulla blockchain significa semplicisticamente eseguire le tre operazioni
		\begin{enumerate}[label*=\arabic*.]
			\item \textbf{Generazione di una chiave privata casuale}
			\item \textbf{Conversione della chiave privata generata in (1) in una chiave pubblica}: un protocollo comune per la conversione
			in questa industria è l'algoritmo ECDSA
			\item \textbf{Conversione della chiave pubblica generata in (2) in un indirizzo di valuta}: protocolli comuni per eseguire questa
			conversione sono le funzioni hash SHA-256, Base58 encoding, Base32 encoding
		\end{enumerate}
	A questo punto qualsiasi valore inviato all'indirizzo generato in (3) è garantito dalla doppia spesa nella blockchain scelta ed accessibile
	solo dal possessore della chiave privata generata in (1).
	\item \textsc{dove}: due fattori critici che interessando l'industria delle valute virtuali sono la concorrenza bancaria ed il 
	banning a livello statale. Anche se l'armonizzazione delle leggi europee in materia sarà di beneficio all'industria sia in termini di 
	tassazione sia in termini di garanzia per l'utilizzatore finale, al momento questo ancora non avviene. Ogni Paese ha la propria
	legislazione in materia ed i propri regimi di tassazione per tutte le operazioni di cambia valuta ed il banning statale va da livelli
	di permissività totale come nell'Unione Europea al banning totale con imprigionamento come in Bangladesh \cite{bitcoinLegality}. 
	\item \textsc{dove}: i paesi asiatici della Corea del Sud, Cina e Giappone sono sempre stati leader nell'industria delle valute virtuali
	negli ultimi 9 anni con un approccio proattivo in materia di tassazione. All'inizio del 2017 in Giappone il bitcoin è stato 
	dichiarato legal tender tuttavia in Cina recentemente sono state dichiarate illegali tutte le operazioni di cambia valuta virtuale ed ICO 
	e gli exchange nel Paese asiatico stanno cessando le loro attività di cambio. 
	\item \textsc{dove}: qualsiasi valutazione del mercato locale dovrebbe comprendere: numero di potenziali utenti, tipo di exchange da 
	incorporare (FIAT <-> CRYPTO oppure CRYPTO <-> CRYPTO), tipi di protocolli da integrare (POW, DPOS, Masternode, altri...), 
	tipi di servizi da offrire (trading, trading avanzato, processore di pagamento, altri...), se l'exchange incorporato è di tipologia
	FIAT <-> CRYPTO numero e tipologia di processori di pagamento accettati (PayPal, OKPAy, MoneyPolo, altri...), numero di exchange
	simili nella regione di incorporazione. 
	\item \textsc{chi}: ci sono numerose opzioni per offrire una polizza assicurativa sui fondi degli utilizzatori finali (Customer Protection):
	creare comprensione tra gli utenti che il possesso delle chiavi private dei fondi mette loro in \textbf{posizione di responsabilità} 
	nei confronti di un'eventuale perdita dei fondi, in quanto nessuno può recuperare le chiavi private una volta che sono perse. 
	Creare comprensione negli utenti per non lasciare i fondi negli exchange ("\textit{Be Your own Bank!!}"), lasciandogli scegliere a quale
	exchange del mercato affidarsi.
	\item \textsc{chi}: mentre è decisamente costoso assicurare la valuta e le operazioni di trasmissione del valore, esempi di customer protection
	in questa industria sono: la piattaforma Kraken ha recentemente permesso agli utenti Mt. Gox di richiedere rimborso tramite
	la loro piattaforma, la comunità NEO offre rimborso agli utenti coinvolti nell'hacking di BitGrail, gli utenti Ethereum
	offrono rimborso agli investitori della DAO, altri casi di hacking...
	\item \textsc{Raccomandazioni finali/conformità legislativa}: se intendete incorporare un exchange FIAT <-> CRYPTO dovete focalizzarvi
	in prima istanza/inanzitutto nella conformità legislativa in materia di AML/KYC nel Paese di incorporazione, le fondazioni
	bitcoin locali possono aiutarvi in focalizzare le vostre operazioni di cambio basate sugli interessi degli utilizzatori finali 
	nel Paese di incorporazione, promuovendo ambienti \textit{cryptocurrency friendly} nell'area di interesse.
\end{enumerate}


%----------------------------------------------------------------------------------------
%	CHAPTER 3: Ripa Exchange
%----------------------------------------------------------------------------------------

\chapterimage{chapter_head_3_Singapore.jpg} % Chapter heading image

\chapter{L'Exchange Ripa}
Ripa Exchange è un crypto asset marketplace implementato seguendo i più avanzati standard industriali seguendo il principio
"\textit{open source, sicuro ed efficiente}". Ripa Exchange si prefigge di servire come piattaforma sicura, UI responsive, 
personalizzabile, facile da usare per gli imprenditori in questo settore abbracciando i principi Open Source e la fiducia pubblica.\\\\
Ripa Exchange è implementato utilizzando il framework Ruby on Rails ed altre tecnologie all'avanguardia e sarà migrato verso
un exchange ibrido-decentralizzato dove tutti gli exchange nella rete Ripa condivideranno la stessa liquidità grazie alla tecnologia R.L.S.P..

\section{Mission}
\begin{quotation}
	``\textit{La nostra mission è di costruire il crypto asset marketplace migliore al mondo con un trading engine ad alte prestazioni 
	e sicuro abbastanza da poter essere utilizzato dagli utenti in tutta fiducia. In aggiunta vogliamo far fare un passo avanti 
	alla tecnologia degli exchange di criptomonete offrendo supporto tecnico e nuove funzionalità. Aiutiamo persone da tutto il mondo
	alla costruzione dei loro exchange personalizzati.}''
\end{quotation}

Il supporto tecnico è sempre apprezato: sentitevi liberi di inoltrare pull-request oppure aprire issue nei nostri repository GitHub.

\section{Funzionalità}
Un libero, trasparente ed internazionalizzato exchange di cryptomonete open source.\\

\tcbset{featureBox/.style={colback=yellow!10!white,colframe=white!20!black,
equal height group=nobefaf,width=(\linewidth-1pt)/3, height=4.5cm, nobeforeafter,
	center title,
	valign=top, halign=left}}
\begin{tcolorbox}[featureBox,
	title=\textsc{Open Source}]

	\small	Tutto il codice sorgente prodotto è rilasciato sotto licenza MIT.\\\vspace{5mm}
	\tiny Ripa Exchange è un'architettura di cambia valuta virtuale personalizzabile che connette semplicemente procedure
	di autentificazione AML/KYC, report ETL ed altri servizi.
\end{tcolorbox}
\begin{tcolorbox}[featureBox,
	title=\textsc{Conforme}]

	\small Standard internazionali AML/KYC.\\\vspace{5mm}
	\tiny La KYC di Ripa Exchange permette procedure KYC semplici a livello bancario e conformi ai requisiti di 
	Customer Due Diligence (CDD).
\end{tcolorbox}
\begin{tcolorbox}[featureBox,
	title=\textsc{Trasparente e Configurabile}]

	\small Personalizza a tuo modo\\\vspace{5mm}
	\tiny Le principali funzionalità sono state racchiuse nel codice sorgente: registrazione pulita ed 
	interfaccia di login semplice, procedure di deposito e prelievo personalizzate, miglior accoppiamento di ordini 
	di vendita ed acquisto, altro... . Queste funzionalità sono già integrate nell'exchange finale e sono pronte all'uso
	senza lavoro aggiuntivo. 
\end{tcolorbox}

\begin{tcolorbox}[featureBox,
	title=\textsc{Internazionalizzazione}]

	\small	Gli utenti possono visualizzare Ripa Exchange nella lingua preferita.\\\vspace{5mm}
	\tiny Supportando diversi linguaggi comuni, Ripa Exchange rende facili le operazioni di trading degli utenti
	nella loro lingua madre. Vi incoraggiamo a contribuire alla varietà linguistica di Ripa Exchange: la comunità
	di utenti trarra beneficio da questa.
\end{tcolorbox}
\begin{tcolorbox}[featureBox,
	title=\textsc{Prova di Solvibilità}]

	\small	PoS facilmente rilasciabile.\\\vspace{5mm}
	\tiny La Prova di Solvibilità (PoS) di Ripa Exchange permette agli utenti di verificare la solvibilità dell'exchange
	senza compromettere la privacy del singolo utente.
\end{tcolorbox}
\begin{tcolorbox}[featureBox,
	title=\textsc{Multi-Accounts Trading}]

	\small	Facilità di configurazione delle valute.\\\vspace{5mm}
	\tiny Ripa Exchange permette di creare molteplici account e scambiare in diverse valute: quest operazione è
	di facile esecuzione in Ripa Exchange.
\end{tcolorbox}

\begin{tcolorbox}[featureBox,
	title=\textsc{Multi-Accounts Utenti}]

	\small	Facilità di configurazione degli account utente.\\\vspace{5mm}
	\tiny Ripa Exchange offre molteplici tipologie di login utente tramite Google, Facebook, Twitter ed segue
	gli standard FIDO Alliance per la sicurezza degli account.
\end{tcolorbox}
\begin{tcolorbox}[featureBox,
	title=\textsc{Exchange Enterprise}]

	\small Avvio piccolo, crescita grande.\\\vspace{5mm}
	\tiny Ripa Exchange offre funzionalità enterprise incluso un matching engine ad alte prestazioni, 
	thread di lavoro distribuito scalabili ed autentificazione a due fattori tramite SMS.
\end{tcolorbox}
\begin{tcolorbox}[featureBox,
	title=\textsc{Funzionale ed Intuitivo}]

	\small	Per il trader novizio, per il trader esperto.\\\vspace{5mm}
	\tiny Interfaccie di registrazione, login e trading pulite. Procedure di deposito e prelievo personalizzate 
	e prova di solvibilità integrata.
\end{tcolorbox}

\section{Analisi Funzionale}
\begin{enumerate}
	\item \textsc{Cosa serve Redis-RabbitMQ}:
	\item \textsc{Cosa serve NodeJS}: 
	\item \textsc{Cosa serve Pusher}:
	\item \textsc{ER MySQL}:
	\item \textsc{Ruby gerarchia directory}:
\end{enumerate}

\section{Stack Tecnologico}
\begin{description}
	\item[\textsc{Ruby on Rails}]:
	\item[\textsc{MySQL}]: 
	\item[\textsc{Redis}]:
	\item[\textsc{RabbitMQ}]:
	\item[\textsc{NodeJS}]:
	\item[\textsc{Pusher}]:
\end{description}

\section{Interfaccia Utente}
La user interface di Ripa Exchange è basata sulla user interface di Peatio una UI responsiva sviluppata in Ruby on Rails e 
completamente separata dall'engine ordini dell'exchange.
Il design per una UI personalizzata sono in corso per offrire agli utenti Ripa Exchange la miglior esperienza utente in ogni
dispositivo: in questa sezione potete trovare alcuni screenshot della UI attuale dell'exchange.

\subsection{Interfaccia End-User}
Di seguito alcuni screenshot dell'interfaccia utente end-user dell'exchange Ripa:
\textcolor{darkred}{questi screenshot sono un lavoro in corso e non è detto che rappresentino l'interfaccia utente del prodotto finale}\\
\begin{center}
	\includegraphics[width=\textwidth*2/3,height=\textheight,keepaspectratio]{RE_TradingUIC}
	\captionof{figure}{UI di trading di Ripa Exchange}
	\includegraphics[width=\textwidth*2/3,height=\textheight,keepaspectratio]{RE_depositBTCC}
	\includegraphics[width=\textwidth*2/3,height=\textheight,keepaspectratio]{RE_historyC}
	\includegraphics[width=\textwidth*2/3,height=\textheight,keepaspectratio]{RE_solvencyC}
	\captionof{figure}{UI di deposio/prelievo, storico ordini e solvibilità}
\end{center}

\subsection{Admin Interface}
Di seguito alcuni screenshot dell'interfaccia utente dell'area amministrativa dell'exchange Ripa:
\textcolor{darkred}{questi screenshot sono un lavoro in corso e non è detto che rappresentino l'interfaccia utente del prodotto finale}\\
\begin{center}
	\includegraphics[width=\textwidth*2/3,height=\textheight,keepaspectratio]{RE_adminDashboardC}
	\includegraphics[width=\textwidth*2/3,height=\textheight,keepaspectratio]{RE_adminDepositsC}
	\includegraphics[width=\textwidth*2/3,height=\textheight,keepaspectratio]{RE_adminWithdrawsC}
	\captionof{figure}{Admin console dashboard, sezioni di deposito/prelievo}
	\includegraphics[width=\textwidth*2/3,height=\textheight,keepaspectratio]{RE_adminUserInfo1C}
	\includegraphics[width=\textwidth*2/3,height=\textheight,keepaspectratio]{RE_adminVerifyAccountC}
	\captionof{figure}{Sezione informazioni utente e verifica utente}
\end{center}

\section{Funzionalità al Lancio}
\begin{description}
	\item[\textsc{CRYPTO <-> CRYPTO}]:
	\item[\textsc{OAuth}]: Facebook, Google, Twitter
	\item[\textsc{FIDO}]: login con standard FIDO Alliance
	\item[\textsc{Protocolli}]: POW, DPOS, Masternode, tether, ERC20
	\item[\textsc{Valute}]: BTC, ETH, DOGE, BCH, TUSD, ARK, LISK, SHIFT, RISE, KAPU, OXY, RIPA, promising ERC20 tokens
	\item[\textsc{Mercati Principali}]: BTC, ETH, ARK
	\item[\textsc{Tipo di Ordini}]: market, limit
\end{description}

\subsection{Funzionalità Future}
\begin{description}
	\item[\textsc{E-Wallets}]: OKPay, NETELLER, MoneyPolo, others...
	\item[\textsc{Funzionalità di Trading Avanzate}]: margin trading, stop loss and take profit, 
	\item[\textsc{FIAT <-> CRYPTO}]
	\item[\textsc{Altre funzionalità}]: VISA/MasterCard, merchants tools, P2P Lending 
\end{description}


\section{Verso un Exchange Decentralizzato...}
Ripa Exchange è un exchange centralizzato che sarà convertito in un ibrido-decentralizzato per creare una rete di exchange
che condividono lo stesso codice sorgente e la stessa liquidità in modo tale che tu possa offrire liquitià ai tuoi clienti 
già dal giorno 1 di apertura degli scambi finanziari.\\

Per offrire i PRO di un exchange centralizzato (supporto di livello platinum, scambio FIAT) e di un exchange decentralizzato
(liquidità) senza i CONTRO di entrambi, dovremo eseguire questa passaggio con lo step intermedio di un exchange 
ibrido-decentralizzato e durante la fase 3 del progetto RipaEx (WP 4-6) esguiremo tutte le analisi funzionali e tecniche 
richieste per costruire il prossimo passo su fondamenta solide. 

%----------------------------------------------------------------------------------------
%	PART
%----------------------------------------------------------------------------------------

%\part{Part Two}

%----------------------------------------------------------------------------------------
%	CHAPTER 4: Ripa Blockchain
%----------------------------------------------------------------------------------------

\chapterimage{chapter_head_4_Dubai.jpg} % Chapter heading image
\chapter{La Blockchain Ripa}
Il progetto RipaEx avrà la propria blockchain che sarà denominata Ripa, securizzata tramite il protocollo DPOS e
con il proprio token \PHP (simbolo XPX) scambiato su di essa che servirà a svolgere le seguenti funzioni:
	\begin{enumerate}
		\item per listare nuove valute nell'exchange
		\item per pubblicizzare nuovi progetti
		\item per comprare gadget RipaEx nello Store Ripa Exchange
		\item per pagare per l'acquisto di beni e servizi in rivenditori autorizzati tramite il RipaEx POV (Punto di Vendita)
		\item per condividere liquidità tra tutti gli exchnage della rete Ripa
	\end{enumerate}

Il token \PHP o XPX è una criptovaluta derivata da ARK, Lisk, Crypti e BitShares con differenze uniche e miglioramenti per 
raggiungere l'obbiettivo di condividere liquidità tra tutti gli exchange della rete Ripa. Tuttavia questo codice eredita
le transazioni semplificate tra le blockchain basate su ARK ed il consenso DPOS. Questo codice sorgente omogeneo 
permette lo sviluppo di potenziali servizi intra-blockchain nella forma di app ARK ed assieme ad altri sistemi addizionali
forniti dagli amministratori delle blockchain rispettive.

\vspace{5mm}
\textsc{\textbf{La blockchain Ripa è un fork di ARK e la creazione di tale blockchain serve per completare l'ecosistema Ripa
permettendo ogni exchange nella rete Ripa di condividere la stessa liquidità. Noi ci fideremo sempre di ARK come nostro
fornitore di tecnologia blockchain per mergiare il loro codice nel nostro per tutto ciò che riguarda la blockchain Ripa}}

\vspace{5mm}

Spiegato l'uso della blockchain Ripa all'interno dell'ecosistema RipaEx e spiegato il legame di partnership tecnologica
RIPA - ARK nelle seguenti sezioni potete trovare le specifiche della blokchain Ripa che sono derivata da ARK ed alcune sviluppate 
"ad hoc" per l'ecosistema Ripa.

\section{Tecnologia Delegated Proof of Stake}
Ripa Blockchain, will inherits the Delegated Proof of Stake (DPoS) consensus system that was first introduced
by BitShares. This consensus algorithm was designed to eliminate the issues associated with Proof of Work (PoW), 
namely the centralization of computing power and the exponentially increasing waste of real world energy. While not
completely decentralized as it relies on consensus by a fixed number of elected delegates, it guarantees a better decentralization 
than Bitcoin. The consensus algorithm implementation is improved over time, evolving into an optimal consensus system.\\

The technical description of the Ripa blockchain is as follow:
\begin{enumerate}
	\item DPoS (Delegated Proof of Stake)
	\begin{description}
		\item - 101 active forging Delegates
		\item - Delegates selected by vote mechanism built into DPoS
		\item - 115,000,000 RIPA - Seeded Genesis Block
	\end{description}
	\item Multi-signature accounts
	\item Constant block reward
	\begin{description}
		\item - 2 \PHP per block
		\item - Inflation Rate (with 8s block times)
		\begin{description}
			\item 6.31\% for the first year
			\item 5.93\% the 2nd year
			\item 4.02\% the 10th year
		\end{description}
		\item - 8-second block time
		\begin{description}
			\item - Decreased block time possible with future upgrades to the core.
		\end{description}
		\item - 25 transactions per block
		\begin{description}
			\item - Increased via soft fork as needed.
		\end{description}
	\end{description}
	\item Routing tables
	\item SmartBridge data field for custom use and bridging blockchains (ARK Contract Execution Service)
	\item Batch transactions\footnote{Future upgrade: when ARK 2.0 will be released \label{note1}}
	\item Custom transaction fees\footnotemark[\value{footnote}]
	\item Native smart contract execution\footnote{Future upgrade: when ARK Virtual Machine will be released}
\end{enumerate}

\vspace{5mm}
\textsc{\textbf{As said Ripa blockchain is an ARK fork and we will entrust ARK as our blockchain technology provider to merge their 
code into our latest features for what concerning the RIPA Blockchain}}, features like:
\begin{description}
	\item[Scaling network] up to the level of major Credit Card networks with core upgrades
	\begin{description}
		\item - Increasing the number of Forging Delegates
		\item - Increasing the Block Size to include more transactions
		\item - Implementation of pre-approval PBFT block concept testnet [codename:
		TwinChain]
		\item - Routing tables, to minimize hops among nodes when blocks are broadcast
		\item - Include forging with RIPA Uncles.	
	\end{description}
	\item[Two node types] is used to secure the RIPA network:
	\begin{description}
		\item[Relay nodes] - Nodes with full API functionality, acting as a backend for the
		feature rich lite clients. Relay nodes do not collect any transaction fee and do
		not have the ability to Forge RIPA Blocks.
		\item[Forging nodes] - Nodes with reduced API functionality, decreasing the
		exposure to potential DDoS attacks on the RIPA Platform. Forging nodes are
		able to Forge RIPA and receive transaction fees.
	\end{description}
	\item[Official lite client] for network access is be provided shortly before the mainnet
	launch including desktop clients (Windows, MacOS, and Linux) and mobile clients
	(Android and iOS).
	\item[OffLine wallet creation]: the network itself does not use a Graphical User Interface by default. Any RIPA
	account can be created offline and managed at no cost with a single device
	(computer, mobile phone, embedded ARM, IoT).
\end{description}

\section{Hierarchical Deterministic (HD) Wallets (BIP32)}
The structure of the public and private key generation follows the same specification
as Bitcoin. A custom implementation of BIP32 for Hierarchical Deterministic Wallets
is provided to RIPA users.

\section{Fees}
The fee for standard transactions is set at \PHP0.1 but the will be flexible in future releases of Ripa Blockchain.
At mainnet launch, a fee structure is provided out of the box to forging delegates with the following rules:
	\begin{description}
		\item - Transaction \PHP0.1
		\item - Vote \PHP1 (101 votes per transaction)
		\item - Second Signature \PHP1
		\item - Multi Signature \PHP1 per signature + \PHP1 per signing account
		\item - Registering a delegate \PHP25.
	\end{description}
All fees are paid to the forging node which processes the block containing those fees.

\section{Ripa Delegates and Delegate Voting}
Any node running the core blockchain code wishing to become a forging node must
register their account within the RIPA network. The fee for this registration is set to
25 \PHP per delegate account registered.
RIPA incorporates a new DPoS voting system originally envisioned by the Crypti
Founders. The RIPA system fee is 1 \PHP per delegate vote. The voting weight of each
wallet will be split evenly between all delegates voted. For example:
\begin{description}
	\item - If a wallet votes for one delegate, that delegate receives 100\% of the wallets
voting weight.
	\item - If that wallet votes for an additional delegate, the entire vote weight is split
evenly between the two delegates at 50\%.
	\item - By adding a third delegate, the voting weight splits again, and each of the
three delegates receives 33.333\% of the voting weight from that wallet.
\end{description}

The 101 forging nodes with the highest number of votes are eligible to Forge RIPA
blocks. This design eliminates the possibility that any single large RIPA holder or an
organization holding large percentages of RIPA are able to gain control over the
entire network by voting for all of their nodes into forging positions, thus effectively
taking complete control over that DPoS Blockchain. Votes from RIPA Tokens held by
RIPA Crew may be used at RIPA Crew's discretion.

\section{Connecting Blockchains}
\subsection{ARK Bridged Blockchains (SmartBridge Technology)}
The ARK Platform does not provide direct support for sidechains or dapp databases.
Instead, a mechanism to bridge together blockchains is provided via a bridging
function built into ARK Core where any blockchain can send and receive trigger
function notices and informational data through the primary ARK network via
custom developed \textbf{SmartBridge(s)} and \textbf{Encoded Listeners}.
\begin{center}
	\includegraphics[width=\textwidth,height=\textheight,keepaspectratio]{ARKSmartBridge}
	\captionof{figure}{ARK SmartBridge Technology}
\end{center}

\subsection{A.C.E.S. - ARK Contract Execution Service}
ACES is a blockchain interoperability platform that provides simple protocols and tools for building a robust blockchain 
service marketplace.

ACES is composed mainly of the following three components:
\begin{description}
	\item[Listeners] ACES Listeners provide a way for all the different blockchain transaction events to be easily 
	consumed via a common REST-ful API. The API allows consumers to create subscriptions and receive blockchain events 
	in real-time using Webhook callbacks.
	\item[Services] ACES Services create and excute Service Contracts, which can be anything from uploading 
	a file to a storage blockchain, performing value transfers, creating smart contracts, executing code on 
	blockchain based computing platforms, or interacting with IoT hardware.
	\item[Marketplace Console] The ACES Marketplace Console is a consumer dashboard for searching and executing 
	service contracts listed on the Marketplace. ACES Service providers can list their service nodes using the Marketplace API.
\end{description}

\begin{center}
	\includegraphics[width=\textwidth,height=\textheight,keepaspectratio]{ACES}
	\captionof{figure}{ARK-BTC A.C.E.S. SmartBridge Implementation Overview}
\end{center}

\section{Ripa Liquidity Service Provider (R.L.S.P)}
Using the SmartBridge technology RipaEx will build a mechanism to share liquidity between all the exchanges in the Ripa network
by writing the single exchange orderbook in the Ripa Blockchain and by executing order matching between all exchanges in the network.

In this way you can have the benefits of a decentralized orderbook (like liquidity) with the benefits of a centralized exchange (like
platinum customer support and FIAT exchange).

\section{Ripa Community Fund}
To permits the born of new exchanges in the Ripa network a Ripa Community Fund - RCF - is created with the following characteristics:
\begin{description}
	\item[Starting Principal]: the RCF will have a starting operating capital of 5\% (5,750,000 \PHP) of the genesis block
	\item[Recurring Participation]: each delegate will contribute to the RCF with 5\% of its forged XPXs
\end{description}

To obtain funds from the RCF to start your own Ripa Exchange you will submit your proposal in the official RCF section in the Ripa forum at
anytime after the first Ripa Exchange instance will be operative.

%----------------------------------------------------------------------------------------
%	CHAPTER 5: Token Sale
%----------------------------------------------------------------------------------------

\chapterimage{chapter_head_5_Tokyo.jpg} % Chapter heading image
\chapter{Token Sale}
\section{Introduction}
The RipaEx XPX token sale will be separated into two phases:
	\begin{description}
		\item[\textsc{PreSale}]: that will run from April to September 2018
		\item[\textsc{RIPA TEC}]: that will run from October to December 2018
	\end{description}

You can easily refer to both phases presale and RIPA TEC as \textit{RipaEx ICO}.

\subsection{Exchange Rate \& Accepted Cryptocurrencies}
During the token sale the exchange rate for the XPX tokens will be \textbf{\PHP/\euro0.10}.\\

The following cryptocurrencies are accepted: \textbf{BTC, ETH, ARK, LISK}.\\

If you want to send us a different cryptocurrency send us your choice and we will tell you immediately if we can make 
the exchange and how to execute the operation in detail (receiving address, exchange rate, quantity)\\

\textsc{Exchange rate is Bitstamp last and is valid 60 minutes from quote}.

\section{How to Invest}
\subsection{XPX Ripa Token PreSale}
Private sale of XPX tokens has \textbf{started on Monday 04/02/2018 at 00:00 GMT and will end on Sunday 09/30/2018 at 23:59 GMT}.

You can buy directly here: 
\begin{center}
	\href{https://www.ripaex.io/privsale}{\textsc{www.ripaex.io/privsale}}
\end{center}

or send us your buying enquiries on our social Slack, Telegram, Facebook, Twitter or on other social networks where we are present.
For all the duration of the private sale period every exchange operation \textbf{will receive a 100\% bonus on its sending amount
until the first 5.000.000 XPX tokens} have been exchanged and will receive a 50\% bonus thereafter.

\subsection{RIPA TEC}
Fundraising through the RIPA TEC platform will follow the schedule:
\begin{description}
	\item[from 10/01/2018 00:00 to 10/14/2018 23:59]: with bonus of 50\% 
	\item[from 10/15/2018 00:00 to 10/28/2018 23:59]: with bonus of 40\% 
	\item[from 10/29/2018 00:00 to 11/11/2018 23:59]: with bonus of 30\% 
	\item[from 11/12/2018 00:00 to 11/25/2018 23:59]: with bonus of 20\% 
	\item[from 11/26/2018 00:00 to 12/09/2018 23:59]: with bonus of 10\% 
	\item[from 12/10/2018 00:00 to 12/23/2018 23:59]: with bonus of 5\% 
\end{description}
Coins distribution will be available at the end of each trading day and performed automatically through the RIPA TEC platform. Minimum
threshold to reach is 25 BTC.\\

\textsc{As the XPX exchange rate is calculated against the euro currency directly you are guaranteed to not loose any invested fund
during the sale period as the XPX tokens will be listed on our exchanges at the beginning of 2019 directly with the same \textbf{\PHP/\euro0.10}
exchange rate agreed during the token sale-ICO phase}.

\section{XPX Ripa Token Distribution}
\textbf{115,000,000} XPX tokens are seeded into the genesis block: distribution of the XPX tokens generated are described into
figure \ref{fig:distribution} and table \ref{tab:distribution}.

\vspace{5mm}
%\begin{center}
	\begin{tikzpicture}
		\pie [rotate = 180, explode={0.2, 0.1, 0.1, 0.1, 0.1, 0.1}, radius = 3]
		{65/PreSale \& RIPA TEC,
		15/RIPA Founders Team, 
		6/for marketing,
		5/for bounties,
		5/for RCF,
		4/ARK Team}
	\end{tikzpicture}
	\captionof{figure}{Ripa token distribution - Division of funds}	
	\label{fig:distribution}
%\end{center}

\vspace{5mm}
\begin{table}[H]
	\centering
	\begin{tabular}{l l l}
		\toprule
		\textbf{Percentage (\%)} & \textbf{Quantity (\PHP)} & \textbf{Purpose} \\
		\midrule
		65		& 74,750,000	& to distribute in PreSale and RIPA TEC	\\
		15      & 17,250,000	& to the RIPA Founders Team	\\
		6       & 6,900,000		& for marketing	\\
		5       & 5,750,000 	& for bounties	\\
		5       & 5,750,000		& for Ripa Community Fund - RCF	\\
		4       & 4,600,000		& to the ARK Team	\\
		\bottomrule
	\end{tabular}
	\captionof{table}{Ripa token distribution}
	\label{tab:distribution}
\end{table}

\vspace{5mm}
\textbf{Tokens not sold during the phases of PreSale and RIPA TEC will be burnt} forever to permits and equal distributions 
of the funds and avoid speculations on the token remaining.

\section{Division of Funds}
The funds collected during the PreSale and RIPA TEC phases will be divided as explained in figure \ref{fig:division} and 
table \ref{tab:division} below

\vspace{5mm}
\begin{tikzpicture}
	\pie [rotate = 180, explode={0.2, 0.1, 0.1}, radius = 3]
	{60/for project development,
	20/for marketing, 
	20/for legal support}
\end{tikzpicture}
\captionof{figure}{Ripa division of funds}	
\label{fig:division}

\vspace{5mm}
\begin{table}[H]
	\centering
	\begin{tabular}{l l}
		\toprule
		\textbf{Percentage (\%)} & \textbf{Purpose} \\
		\midrule
		60		& for project development	\\
		20		& for marketing	\\
		20		& for legal support	\\
		\bottomrule
	\end{tabular}
	\captionof{table}{Division of funds}
	\label{tab:division}
\end{table}

\vspace{5mm}
where the row project development funds allocation is described in table \ref{tab:focus} below.

\vspace{5mm}
\begin{table}[H]
	\centering
	\begin{tabular}{l l}
		\toprule
		\textbf{Percentage (\%)} & \textbf{Purpose} \\
		\midrule
		50		& functional analysis, technical analysis, development, testing	\\
		20		& technical support	\\
		15		& infrastructure	\\
		10		& security	\\
		5		& analytics	\\
		\bottomrule
	\end{tabular}
	\captionof{table}{Project development focus}
	\label{tab:focus}
\end{table}

%----------------------------------------------------------------------------------------
%	CHAPTER 6: Economic Projections
%----------------------------------------------------------------------------------------

\chapterimage{chapter_head_6_Frankfurt.jpg} % Chapter heading image

\chapter{Business Model}

\section{Project Milestones}
Following the matrix of development per area of interest based on the funds collected during the phases of PreSale and RIPA TEC:\\
\begin{tcbraster}[raster columns=5,raster rows=1,raster height=0.8cm,
	valign=center, halign=center,
	enhanced,size=small,sharp corners,colframe=azure(colorwheel),coltext=white,
	colback=azure(colorwheel),fit algorithm=hybrid* ]
	\tcboxfit{}
	\tcboxfit{\textbf{25 BTC}}
	\tcboxfit{\textbf{250 BTC}}
	\tcboxfit{\textbf{1000 BTC}}
	\tcboxfit{\textbf{2000 BTC}}
\end{tcbraster}
\begin{tcbraster}[raster columns=5,raster rows=4,raster height=10cm,
	valign=center, halign=center,
	enhanced,size=small,sharp corners,colframe=silver,coltext=black,
	colback=silver,fit algorithm=hybrid* ]
	\tcboxfit{\textsc{\textbf{Development}}}
	\tcboxfit{\tcbfontsize{0.8} CRYPTO <-> CRYPTO exchange release with 25 cryptocurrencies supported and 3 main trading pairs}
	\tcboxfit{\tcbfontsize{0.8} CRYPTO <-> FIAT exchange release with MasterCard prepaid card. 50 cryptocurrencies supported and 3 main trading pairs}
	\tcboxfit{\tcbfontsize{0.8} Advanced trading features. Open of a second and third Ripa Exchanges to join the Ripa network}
	\tcboxfit{\tcbfontsize{0.8} Open of ten Ripa Exchanges across the Globe}

	\tcboxfit{\textsc{\textbf{Marketing}}}
	\tcboxfit{\tcbfontsize{0.8} Support from one agency, adv campaigns on targeted channels}
	\tcboxfit{\tcbfontsize{0.8} Support from three agencies, adv capaigns pushing harder, store with RipaEx gadgets}
	\tcboxfit{\tcbfontsize{0.8} International presentation event in London}
	\tcboxfit{\tcbfontsize{0.8} Support from five agencies, targeting all UN countries, presentation event in Tokyo}

	\tcboxfit{\textsc{\textbf{Legal}}}
	\tcboxfit{\tcbfontsize{0.8} Legal department working on AML/KYC international standards}
	\tcboxfit{\tcbfontsize{0.8} Legal departments working on AML/KYC international standards in New York, London, Hong Kong}
	\tcboxfit{\tcbfontsize{0.8} International Government Agencies cooperation}
	\tcboxfit{\tcbfontsize{0.8} Banking parnership for defining AML/KYC standards}

	\tcboxfit{\textsc{\textbf{Blockchain}}}
	\tcboxfit{\tcbfontsize{0.8} DevNET development for RLSP and future ARK releases}
	\tcboxfit{\tcbfontsize{0.8} Contributions to ARK development for AVM}
	\tcboxfit{\tcbfontsize{0.8} Contributions to ARK development for standard timing releases}
	\tcboxfit{\tcbfontsize{0.8} Partnership with ARK for blockchain technology developments}
\end{tcbraster}

\section{Market Overview}

\section{Local Market Analysis}

\section{5 Years Projections}

\section{XPX Token Economics}

%----------------------------------------------------------------------------------------
%	CHAPTER 7: Conclusion
%----------------------------------------------------------------------------------------

\chapterimage{chapter_head_7_LosAngeles.jpg} % Chapter heading image

\chapter{Team \& Conclusion}

\section{Team}
\begin{tcbraster}[raster columns=1,raster rows=1,raster height=0.4cm,
	valign=center, halign=center,
	enhanced,size=small,sharp corners,colframe=azure(colorwheel),coltext=white,
	colback=azure(colorwheel),fit algorithm=hybrid* ]
	\tcboxfit{}
\end{tcbraster}
\begin{tcbraster}[raster columns=2,raster rows=4,raster height=10cm,
	valign=center, halign=left,
	enhanced,size=small,sharp corners,colframe=silver,coltext=black,
	colback=silver,fit algorithm=hybrid* ]
	\tcboxfit{\tcbfontsize{0.8} \textbf{Giovanni Silvestri @ BitNow}\\
	Founder and Head of the Project\\
	10+ Years experience in IT Field,\\
	4+ Years in Cryptocurrencies Consultancy and Private Exchanging,\\
	2+ Years Experience in finance\\
	Location: Treviso (TV), Italy\\
	Bitcointalk
	Linkedin}
	\tcboxfit{\tcbfontsize{0.8} CRYPTO <-> FIAT exchange release with MasterCard prepaid card. 50 cryptocurrencies supported and 3 main trading pairs}

	\tcboxfit{\tcbfontsize{0.8} Support from three agencies, adv capaigns pushing harder, store with RipaEx gadgets}
	\tcboxfit{\tcbfontsize{0.8} International presentation event in London}

	\tcboxfit{\tcbfontsize{0.8} International Government Agencies cooperation}
	\tcboxfit{\tcbfontsize{0.8} Banking parnership for defining AML/KYC standards}

	\tcboxfit{\tcbfontsize{0.8} Contributions to ARK development for standard timing releases}
	\tcboxfit{\tcbfontsize{0.8} Partnership with ARK for blockchain technology developments}
\end{tcbraster}


\section{Recommendations}

\section{Social}

%----------------------------------------------------------------------------------------
%	CHAPTER 8: Legal
%----------------------------------------------------------------------------------------

\chapterimage{chapter_head_8_KualaLumpur.jpg} % Chapter heading image

\chapter{Legal}

%----------------------------------------------------------------------------------------
%	BIBLIOGRAPHY
%----------------------------------------------------------------------------------------

\chapterimage{chapter_head_9_Milan.jpg} % Chapter heading image

\addcontentsline{toc}{chapter}{\textcolor{trolleygrey}{References}}
\nocite{*}
\printbibliography[title={References}]
%------------------------------------------------


%----------------------------------------------------------------------------------------
%	COPYRIGHT PAGE
%----------------------------------------------------------------------------------------

\newpage
~\vfill
\thispagestyle{empty}

\noindent Copyright \copyright\ 2018 Ripa Founder Team\\ % Copyright notice

\noindent \textsc{Pubblicato da Ripa Founder Team}\\ % Publisher

\noindent \textsc{www.ripaex.io}\\ % URL

\noindent Materiale protetto dalla Licenza MIT (di seguito la ``Licenza''). Non dovrai usare questo file eccetto in conformità con la Licenza. Puoi ottenere una copia della Licenza a \url{https://opensource.org/licenses/MIT}. A meno che non sia richiesto dalla lagislazione applicata o concordato via scritto, il software distribuito con la Licenza è distribuito \textsc{``così com'è'' senza garanzie di alcun tipo}, esplicite o implicite. Fate riferimento alla Licenza per il linguaggio specifico che governa i permessi e le limitazione indicati dalla Licenza.\\ 

\noindent \textit{Prima pubblicazione, Marzo 2018} % Printing/edition date

\end{document}
